% SEGMENT OLP-0078-S01
% Part: first-order-logic
% Chapter: sequent-calculus
% Section: provability-consistency

\documentclass[../../../include/open-logic-section]{subfiles}

\begin{document}

\iftag{FOL}
      {\olfileid{fol}{seq}{prv}}
      {\olfileid{pl}{seq}{prv}}

\olsection{\usetoken{S}{derivability} மற்றும் முரணின்மை}

!!{derivability} உறவின் பல பண்புகளை இப்போது நிறுவுவோம். அவை தனித்தும்
சுவையானவை; மேலும் அவை ஒவ்வொன்றும் நிறைவுத் தேற்றத்தின் நிறுவலில் ஒரு பங்கு
வகிக்கும்.

% SEGMENT OLP-0078-S02
\begin{prop}\ollabel{prop:provability-contr}
  $\Gamma \Proves !A$ என்றும், $\Gamma \cup \{!A\}$ முரணுடையது என்றும்
  இருந்தால், $\Gamma$ முரணுடையது.
\end{prop}

\begin{proof}
$\Log{LK}$ ஆனது $\Gamma_0 \Sequent !A$, $!A, \Gamma_1 \Sequent \quad$
ஆகியவற்றை !!{derive}s எனுமாறு முடிவுறு $\Gamma_0$, $\Gamma_1 \subseteq
\Gamma$ உள்ளன. $\Gamma_0 \Sequent !A$ இன் $\Log{LK}$-!!{derivation}[acc]
~$\pi_0$ என்றும், $\Gamma_1, !A \Sequent \quad$ இன்
$\Log{LK}$-!!{derivation}[acc]~$\pi_1$ என்றும் கொள்க. அப்போது பின்வருமாறு
!!{derive} இயலும்:
\begin{prooftree}
\AxiomC{}
\RightLabel{$\pi_0$}
\Deduce$ \Gamma_0 \fCenter !A $
\AxiomC{}
\RightLabel{$\pi_1$}
\Deduce$!A, \Gamma_1 \fCenter $
\RightLabel{\Cut}
\BinaryInf$ \Gamma_0,\Gamma_1 \fCenter $
\end{prooftree}

$\Gamma_0 \subseteq \Gamma$, $\Gamma_1 \subseteq \Gamma$ என்பதால்,
$\Gamma_0 \cup \Gamma_1 \subseteq \Gamma$; ஆகவே $\Gamma$ முரணுடையது.
\end{proof}

% SEGMENT OLP-0078-S03
\begin{prop}
\ollabel{prop:prov-incons}
$\Gamma \Proves !A$ என்பது, அப்போதும் அப்போது மட்டுமே, $\Gamma \cup
\{\lnot !A\}$ முரணுடையது.
\end{prop}

\begin{proof}
முதலில் $\Gamma \Proves !A$ என்க; அதாவது, $\Gamma \Sequent !A$ இன்
!!a{derivation}~$\pi_0$ உள்ளது. ஒரு $\LeftR{\lnot}$ விதியைச் சேர்ப்பதன்
மூலம்~$\lnot !A, \Gamma \Sequent \quad$ இன் !!a{derivation} ஒன்றைப்
பெறுகிறோம்; அதாவது, $\Gamma \cup \{\lnot !A\}$ முரணுடையது.

$\Gamma \cup \{\lnot !A\}$ முரணுடையது எனில், $\lnot !A, \Gamma
\Sequent \quad$ இன் !!a{derivation}~$\pi_1$ உள்ளது. பின்வருவது $\Gamma
\Sequent !A$ இன் !!a{derivation}:
\begin{prooftree}
  \Axiom$!A \fCenter !A$
  \RightLabel{\RightR{\lnot}}
  \UnaryInf$\fCenter !A, \lnot !A$
  \AxiomC{}
  \RightLabel{$\pi_1$}
  \Deduce$\lnot !A, \Gamma \fCenter$
  \RightLabel{\Cut}
  \BinaryInf$\Gamma \fCenter !A$
\end{prooftree}
\end{proof}

% SEGMENT OLP-0078-S04
\begin{prob}
$\Gamma \Proves \lnot !A$ என்பது, அப்போதும் அப்போது மட்டுமே, $\Gamma
\cup \{!A\}$ முரணுடையது என்பதை நிறுவுக.
\end{prob}

% SEGMENT OLP-0078-S05
\begin{prop}\ollabel{prop:explicit-inc}
  $\Gamma \Proves !A$ என்றும் $\lnot !A \in \Gamma$ என்றும் இருந்தால்,
  $\Gamma$ முரணுடையது.
\end{prop}

\begin{proof}
  $\Gamma \Proves !A$ என்றும் $\lnot !A \in \Gamma$ என்றும் கொள்க.
  $\Gamma_0 \Sequent !A$ என்ற தொடரணியின் !!a{derivation}~$\pi$ ஒன்று
  உள்ளது. $\lnot !A, \Gamma_0 \Sequent \quad$ தொடரணியும்
  !!{derivable}:
  \begin{prooftree}
    \AxiomC{}
    \RightLabel{$\pi$}
    \Deduce$\Gamma_0 \fCenter !A$
    \Axiom$!A \fCenter !A$
    \RightLabel{\LeftR{\lnot}}
    \UnaryInf$\lnot !A, !A \fCenter$
    \RightLabel{\LeftR{\Exchange}}
    \UnaryInf$!A, \lnot !A \fCenter$
    \RightLabel{\Cut}
    \BinaryInf$\Gamma_0, \lnot !A \fCenter$
  \end{prooftree}
  $\lnot !A \in \Gamma$, $\Gamma_0 \subseteq \Gamma$ என்பதால், இது
  $\Gamma$ முரணுடையது என்பதைக் காட்டுகிறது.
\end{proof}

% SEGMENT OLP-0078-S06
\begin{prop}\ollabel{prop:provability-exhaustive}
  $\Gamma \cup \{!A\}$, $\Gamma \cup \{\lnot !A\}$ ஆகிய இரண்டும்
  முரணுடையவை எனில், $\Gamma$ முரணுடையது.
\end{prop}

\begin{proof}
$\Gamma_0 \subseteq \Gamma$, $\Gamma_1 \subseteq \Gamma$ என்ற முடிவுறு
கணங்களும், முறையே $!A, \Gamma_0 \Sequent \quad$, $\lnot !A, \Gamma_1
\Sequent \quad$ ஆகியவற்றின் $\Log{LK}$-!!{derivation}s $\pi_0$, $\pi_1$
உம் உள்ளன. அப்போது பின்வருமாறு !!{derive} இயலும்:
\begin{prooftree}
\AxiomC{}
\RightLabel{$\pi_0$}
\Deduce$ !A, \Gamma_0 \fCenter $
\RightLabel{\RightR{\lnot}}
\UnaryInf$ \Gamma_0 \fCenter \lnot !A$
\AxiomC{}
\RightLabel{$\pi_1$}
\Deduce$\lnot !A, \Gamma_1 \fCenter  $
\RightLabel{\Cut}
\BinaryInf$ \Gamma_0, \Gamma_1 \fCenter $
\end{prooftree}
$\Gamma_0 \subseteq \Gamma$, $\Gamma_1 \subseteq \Gamma$ என்பதால்,
$\Gamma_0 \cup \Gamma_1 \subseteq \Gamma$. ஆகவே $\Gamma$ முரணுடையது.
\end{proof}

\end{document}
